% Contribution owned by /root/singular_fibres.  Finding IDs SN-44-01--06.
% Uses only standard amsmath/amssymb/amsthm commands; bibliography entries
% are in bibliography.tex in this directory.

\section{Cote 44: the corrected height formula and the three nodal terms}
\label{sec:singular-fibres}

The text to be corrected is Grothendieck, cote 44, scan 004, Theorem 5
(manuscript page 2).  HODGE\_SIGN, AR-01, already disproves its displayed
formula under its explicitly transcribed hypotheses.  We do not repeat
that counterexample.  We construct the missing corrections, including a
normalization of the height, and compute them individually.  The
qualification ``geometrically integral fibres'' does not remove these
terms: their carriers can be non-Cartier cycles inside integral fibres.

\subsection{Conventions and the cited height comparison}

Let $k$ be an algebraically closed field of arbitrary characteristic.
Let $f:M\to B$ be a projective morphism over $k$,
where $M$ is smooth projective of dimension $2m$ and $B$ is a smooth
projective connected curve.  Assume that $f$ is flat and its generic
fibre $Y/K$, $K=k(B)$, is smooth and geometrically connected.  All Chow
groups in this subsection have rational coefficients.  An extension
$\alpha\in\operatorname{CH}^{m}(M)_{\mathbb Q}$ of a class
$z\in\operatorname{CH}^{m}(Y)_{\mathbb Q}$ is \emph{admissible} if
\[
 \deg(\alpha v)=0
 \quad\hbox{for every codimension-$m$ cycle $v$ supported on a fibre.}
 \tag{44.1}\label{eq:44-admissible}
\]
For algebraically trivial $z$, such extensions exist.  The degree of the
intersection of two admissible extensions is independent of the
extensions and of the smooth model.  These are the cited results of
Kahn, \emph{Refined height pairing}, Definition 2.2, Proposition 2.8,
Definition 2.10 and (2-9), Propositions 3.6 and 3.8, and Theorem 5.6(a)
\cite{KahnHeight2024}.  We write this \emph{unsigned geometric height} as
$h_{\mathrm{geom}}(z,z')$.  Existence uses algebraic equivalence in the
present application; it does not assume that all homologically trivial
cycles have numerically trivial extensions.

For the next comparison only, assume $k=\mathbb C$.
Here is an independent identification of this height, so that the
correction formula below is not a definition masquerading as a
comparison.  Put $P=\operatorname{Pic}^{m}_{Y}$, the higher Picard variety
for incidence equivalence, and let
$\theta:\operatorname{CH}^{m}_{\mathrm{alg}}(Y)\to P(\overline{\mathbb C(B)})$
be its Abel--Jacobi map.  Choose a Poincare $m$-cycle on $P\times Y$
whose induced endomorphism of $P$ is multiplication by $k_Y^m>0$.
Its transpose induces an isogeny $\lambda_Y^m:P\to P^\vee$.
The rational homomorphism $\lambda_Y^m/k_Y^m$ is the normalized
incidence duality; the integer is necessary when a universal cycle
does not exist with multiplicity one.  With the geometric Poincare
biextension normalization of the bilinear Neron--Tate pairing,
\[
 h_{\mathrm{geom}}(z,z')
 =\frac{1}{k_Y^m}
    \big\langle\theta z,\lambda_Y^m\theta z'\big\rangle_{\mathcal P}.
 \tag{44.2}\label{eq:44-nt-comparison}
\]
This is Chen, \emph{Degeneration of the archimedean height pairing of
algebraically trivial cycles}, arXiv:2512.22788v1, Theorem 7.9
(also Theorem 7.19(a),(b)); see Lemma 3.16, Definition 3.17,
Lemma 3.19 and Remark 7.8 for the normalization and descent
\cite{ChenHeight2025}.  It is the geometric form of Kunnemann's
\emph{Higher Picard varieties and the height pairing}, Theorem 8.2.
The comparison in \eqref{eq:44-nt-comparison} is a cited theorem, not a
new result of this run.  It does \emph{not} use the Griffiths incidence
conjecture: that conjecture enters Chen's further comparison with the
Lear extension of the Hodge-theoretic biextension, Theorem 7.19(c).

We use the sign convention of cote 44 by putting
\[
 \langle z,z'\rangle_G=(-1)^m h_{\mathrm{geom}}(z,z').
 \tag{44.3}\label{eq:44-signed-height}
\]
The pairing $\langle\ ,\ \rangle_{\mathcal P}$ is the degree of the
extended Poincare biextension; there is no additional factor $1/2$.
If a convention for the polarization of a quadratic height includes
that factor, the whole identity, including its right side, is divided
by $2$.  The assertion here is an identity for the canonical incidence
duality.  Its identification with a positive polarization requires the
primitive signs noted by Grothendieck in cote 44, scan 008; no
positivity of every unmodified higher Picard autoduality is assumed.
For curves, $k_Y^1=1$ and $\lambda_Y^1$ is the negative of the usual
principal polarization (Chen, Example 3.20), which checks the sign.
In arbitrary characteristic (44.3) is the signed version of Kahn's
geometric height.  Formula (44.2), as cited here, is a comparison
over complex function fields.  SN-44-06 proves its extension to
arbitrary characteristic separately, by reduction to curves and
with the finite-extension normalization stated there.  In particular
the geometric statement below does not require descent of a higher
Picard variety over an imperfect function field.

\subsection{SN-44-01: the corrected Theorem 5}
\label{sn:44-01}

\paragraph{Statement.}
Let $X/k$ be smooth projective of dimension $2m$, with a pencil
whose smooth base locus has codimension two.  Let
$\pi:M\to X$ be its blowup and $f:M\to\mathbb P^1$ the resulting
morphism.  Assume the generic member $Y$ is smooth and geometrically
connected.  Let $Z,Z'\in\operatorname{CH}^{m}(X)_{\mathbb Q}$ have
geometrically algebraically trivial restrictions $z,z'$ to $Y$.
Here $k$ is algebraically closed of arbitrary characteristic and the
smooth model just specified is part of the hypotheses; no general
resolution-of-singularities assertion in positive characteristic is
being used.  If algebraic triviality is initially witnessed over a
finite extension of $K$, transfer gives the same admissibility
assertions with rational coefficients over $K$ (Kahn, Theorem 5.6(a)
and Proposition 3.9).

Let $V_s\subset N^m(M)_{\mathbb Q}$ be the image of cycles supported
on the fibre $F_s$, where $N^m$ denotes cycles modulo numerical
equivalence on $M$.  Let $R_s$ be the radical of the intersection
pairing on $V_s$.  For each singular fibre choose actual vertical
cycles $D_{s,1},\ldots,D_{s,r_s}$ whose numerical images give a basis
of $V_s/R_s$, and put
\[
 G_s=(\deg D_{s,i}D_{s,j})_{i,j},\qquad
 a_s(Z)=(\deg(\pi^*Z\,D_{s,i}))_i.
\]
An empty matrix contributes zero.  Then
\[
 \boxed{\quad
 \langle z,z'\rangle_G
 =(-1)^m\left(\deg_X(ZZ')
       -\sum_s a_s(Z)^{\mathsf t}G_s^{-1}a_s(Z')\right).
 \quad}                                                     \tag{44.4}
 \label{eq:44-corrected}
\]
The summands and their sum are independent of all the indicated
choices.  This is the corrected geometric Theorem 5 with specified
height and specified extensions.  Formula (44.2) over $\mathbb C$,
and its separately proved extension in SN-44-06 in arbitrary
characteristic, identify its height with the indicated Neron--Tate
pairing, using the full finite-extension degree normalization.
It allows integral singular fibres and does
not require that their vertical cycles be Cartier.

\begin{proof}
The space $N^m(M)_{\mathbb Q}$ is finite dimensional, with dimension
at most the $2m$-th $\ell$-adic Betti number
($\ell\ne\operatorname{char}k$).  To see the bound precisely, take
any finite numerically independent set of rational cycle classes.
The definition of numerical equivalence supplies complementary
cycles whose rational intersection matrix has full rank.  Its
nonzero minor remains nonzero over $\mathbb Q_\ell$, so the cycle
classes are linearly independent over $\mathbb Q_\ell$ in
$H^{2m}(M,\mathbb Q_\ell(m))$.  This bounds every such finite set.
Hence every $V_s$ is finite dimensional and its quotient by the
radical has a nondegenerate symmetric form.
The matrix $G_s$ is therefore invertible.

Two cycles supported on distinct fibres have intersection zero.
For a smooth fibre $i:F\hookrightarrow M$, every vertical cycle is
$i_*u$, and the self-intersection formula gives
\[
 (i_*u)(i_*v)=i_*\bigl(uv\,c_1(N_{F/M})\bigr)=0,
\]
since $N_{F/M}$ is trivial.  Thus all smooth-fibre classes lie in the
radical of the full vertical space $V=\sum_sV_s$.  Moreover
$\bigoplus_s V_s/R_s\to V/\operatorname{rad}V$ is an isomorphism:
surjectivity is immediate, and a relation among the summands pairs
separately with each $V_s$, forcing each term to be radical.  The same
argument identifies the radical of $V$ with the sum of the $R_s$ and
the smooth-fibre spaces.  Only finitely many singular fibres occur.

Set $\alpha=\pi^*Z$, $\beta=\pi^*Z'$.  By the cited existence theorem,
choose an admissible extension $\alpha_0$ of $z$.  Localization for Chow
groups (Fulton, Proposition 1.8) shows that
$\alpha-\alpha_0$ is represented by a finite sum of fibre-supported
cycles.  Consequently $\alpha r=0$ for every $r\in\operatorname{rad}V$:
the $\alpha_0$ term vanishes by admissibility and the difference term
vanishes by the definition of the radical.  This proves in particular
that $a_s(Z)$ does not depend on the choices modulo $R_s$.

For the column $c_s=G_s^{-1}a_s(Z)$ define actual Chow classes
\[
 w_Z=\sum_{s,i}(c_s)_iD_{s,i},\qquad
 \alpha^h=\alpha-w_Z.
\]
They satisfy $\alpha^h D_{s,j}=0$ for every $s,j$, and $\alpha^h r=0$
for every vertical radical class $r$.  Thus $\alpha^h$ is admissible
and still restricts to $z$.  Apply the identical construction to
$\beta$, obtaining $w_{Z'}$ and $\beta^h$.

Orthogonality to $V$ gives
\[
 \deg(\alpha\beta)=\deg(\alpha^h\beta^h)+\deg(w_Zw_{Z'}),\qquad
 \deg(w_Zw_{Z'})=\sum_s a_s(Z)^{\mathsf t}G_s^{-1}a_s(Z').
\]
The first term is $h_{\mathrm{geom}}(z,z')$ by the cited construction;
the projection formula gives $\deg(\alpha\beta)=\deg_X(ZZ')$.
Equation \eqref{eq:44-signed-height} proves
\eqref{eq:44-corrected}; over $\mathbb C$ also apply
\eqref{eq:44-nt-comparison}.  A change of basis conjugates $G_s$
and changes the two columns contragrediently, leaving the scalar
unchanged; replacing a representative by a radical class has already
been handled.  This also proves the asserted choice independence.
\end{proof}

\paragraph{Other codimensions.}
For the hyperplane pencil in the original theorem, in its notation
$i+j=\dim X$, if $i\ne j$, one of $2i,2j$
is strictly less than $\dim X$.  Weak Lefschetz makes restriction in
that degree injective.  Algebraic triviality on $Y$ therefore makes
one original cycle homologically trivial on $X$.  Its pullback to
$M$ is admissible and has zero intersection with every complementary
cycle.  The geometric height and the raw intersection are both zero.
This proves the vanishing recorded in the margin of scan 004, with
the precise comparison above.  The correction issue is concentrated
in the middle case.

\subsection{SN-44-02: the vertical intersection form of an ordinary node}
\label{sn:44-02}

\paragraph{Statement.}
Let $m\ge1$, and let $f:M\to B$ have smooth projective total space of dimension
$2m$ over an algebraically closed field of characteristic different
from $2$.  Suppose $F_s$ is smooth except for one ordinary double
point $p_s$.  Blow up $p_s$ in $M$.  The exceptional divisor is
$E_s\simeq\mathbb P^{2m-1}$, and its intersection with the strict
transform $\overline F_s$ is the smooth quadric
$Q_s\simeq Q^{2m-2}$.  Choose its two middle ruling classes $a,b$.
For $m=1$, $Q_s$ consists of two reduced points and $a,b$ are their
classes; the smooth strict transform $\overline F_s$ may be disconnected.

For an $m$-dimensional cycle $D$ supported on $F_s$, take its strict
transform $\overline D=\overline i_*d$ and write
\[
 j^*d=u a+v b\quad\hbox{in }\operatorname{CH}^{m-1}(Q_s)_{\mathbb Q},
 \qquad \nu_s(D)=u-v.                                      \tag{44.5}
\]
Here $j:Q_s\hookrightarrow\overline F_s$ and
$\overline i:\overline F_s\hookrightarrow\operatorname{Bl}_{p_s}M$.
Then for any two such cycles
\[
 \boxed{\quad\deg_M(DD')=(-1)^m\nu_s(D)\nu_s(D').\quad}       \tag{44.6}
 \label{eq:44-node-form}
\]
The functional $\nu_s$ is well defined on $V_s$ up to simultaneous
change of sign.  Thus $V_s/R_s$ is either zero or one dimensional.
If it is nonzero, a rational cycle $D_s$ can be normalized by
$\nu_s(D_s)=1$, and then $D_s^2=(-1)^m$.

\begin{proof}
Write $p:\widetilde M\to M$ for this blowup, $E$ for its exceptional
divisor, and $h=c_1(\mathcal O_E(1))$.  The Chow group of the quadric
in the indicated degree is freely generated by its two ruling
classes.  Their elementary intersection table is
\[
 \begin{gathered}
 a+b=h^{m-1}|_Q,\qquad
 \deg(a h^{m-1})=\deg(b h^{m-1})=1,\\
 a^2=b^2=\frac{1+(-1)^{m-1}}2,\qquad
 ab=\frac{1-(-1)^{m-1}}2.
 \end{gathered}                                             \tag{44.7}
\]
For $m\ge2$, this is the quadric calculation already recorded and
proved in HODGE\_SIGN, SG-02, applied with its $r=m-1$; it is used
here as a cited calculation.  For $m=1$ the two point classes are
orthogonal idempotents, which gives the same table directly.
In particular each ruling plane has degree one.

Put $\delta_D=u+v$.  The restriction of $\overline D$ to $E$ is
$\delta_Dh^m$, because each ruling plane has degree one in $E$.
The blowup formula for Chow groups says that the kernel of
$p_*:\operatorname{CH}^{m}(\widetilde M)\to\operatorname{CH}^{m}(M)$
is generated by $E^m$.  Therefore $p^*D-\overline D=cE^m$.
Its restriction to $E$ is zero: the restriction of $p^*D$ factors
through the positive-codimension Chow group of a point.  Since
$E|_E=-h$, we obtain
\[
 p^*D=\overline D+(-1)^{m+1}\delta_D E^m.                   \tag{44.8}
\]
These uses of the blowup and self-intersection formulas are precisely
Fulton, Proposition 6.7 and Corollary 6.3; their displayed
specialization here also verifies the sign directly.

The nodal fibre has multiplicity two at $p_s$, so
$p^*F_s=\overline F_s+2E$.  Because a fibre divisor restricts
trivially to a fibre,
\[
 N_{\overline F_s/\widetilde M}=\mathcal O_{\overline F_s}(-2Q).
\]
Consequently
\[
 \deg(\overline D\,\overline D')=-2\deg_Q((ua+vb)(u'a+v'b)).
 \tag{44.9}
\]
Furthermore $E^{2m}=-1$ and
$\deg(\overline D E^m)=(-1)^{m-1}\delta_D$.  On expanding
$\deg(p^*D\,p^*D')$ with formulas (44.8)--(44.9), the two cross terms and
the exceptional square contribute $\delta_D\delta_{D'}$ in total.
Thus, more explicitly,
\[
 \deg(DD')=(u+v)(u'+v')-2\deg_Q((ua+vb)(u'a+v'b))
          =(-1)^m(u-v)(u'-v').
\]
This proves the claimed formula without a local factor $1/2$.

Strict transforms can change under rational equivalence by cycles
supported on $Q$.  For $m\ge2$, their restrictions to $Q$ add
multiples of $h^{m-1}=a+b$, since
$N_{Q/\overline F_s}=\mathcal O_Q(-1)$ and
$\operatorname{CH}^{m-2}(Q)$ is generated by $h^{m-2}$.
For $m=1$ the possible correction group is
$\operatorname{CH}^{-1}(Q)=0$; equivalently, there are no
one-dimensional cycles supported on the two-point scheme $Q$.
Hence $u-v$ is unchanged.  Independently, the formula just proved implies
that a numerically zero vertical class has $\nu_s(D)^2=0$ and hence
$\nu_s(D)=0$; it descends to $V_s$.  The radical is exactly its
kernel whenever its image is nonzero, and otherwise the entire
form is zero.  Scaling any cycle with nonzero $\nu_s$ proves the
last assertion.  Interchanging the two rulings changes all $\nu_s$
by a common sign and leaves every formula unchanged.
\end{proof}

\subsection{SN-44-03: the two vanishing-cycle cases}
\label{sn:44-03}

Assume now that the pencil is a Lefschetz pencil over $\mathbb C$,
with fibres of dimension $2m-1\ge1$.  Let $E_{\mathrm{van}}$ denote
the span of its vanishing cycles in $H^{2m-1}(Y_t,\mathbb Q)$.
This notation avoids confusion with the exceptional divisors in the
previous proof.  Let $S_*$ be the singular fibres for which the
\emph{actual algebraic} vertical space has nonzero $\nu_s$.
Choose the normalized $D_s$ of SN-44-02 and set
\[
 c_s(Z)=\deg(\pi^*Z\,D_s).
\]
Then the corrected Theorem 5 takes the particularly simple form
\[
 \boxed{\quad
 \langle z,z'\rangle_G
   =(-1)^m\deg_X(ZZ')-\sum_{s\in S_*}c_s(Z)c_s(Z').
 \quad}                                                    \tag{44.10}
 \label{eq:44-node-correction}
\]
Its admissible extension is
$\pi^*Z-(-1)^m\sum_{s\in S_*}c_s(Z)D_s$.

If $E_{\mathrm{van}}\ne0$, every local vanishing cycle is nonzero,
$S_*=\varnothing$, and the uncorrected formula is valid with the
height convention (44.3).  If $E_{\mathrm{van}}=0$, the height in
(44.10) is zero, and
\[
 (-1)^m\deg_X(ZZ')=\sum_{s\in S_*}c_s(Z)c_s(Z').             \tag{44.11}
\]
No claim that a prescribed integral generator of a local cohomology
group is algebraic is hidden in the notation: $D_s$ is chosen only
from the actual image of cycles.  A nonzero integral value of
$\nu_s$ may require rational scaling.  In SN-44-04 its value is
exactly one at all three fibres, so no denominator occurs.

\begin{proof}
Equation (44.10) follows from (44.4) and $G_s=((-1)^m)$.
The local specialization sequence at a node is Deligne, \emph{Weil I},
(4.1), or in etale notation (4.3.3).  Its homology dual identifies
the image of the local defect map from middle cycles of $F_s$ with
the kernel of
$\mathbb Q\longrightarrow H_{2m-1}(Y_t,\mathbb Q)$,
$1\mapsto\delta_s$.  Here is the link with the cycle calculation
(44.5).  For $m\ge2$, the circle-bundle sequence for
$\mathcal O_Q(-1)$ identifies the relevant local quotient with
\[
 H^{2m-2}(Q,\mathbb Q)\big/
       hH^{2m-4}(Q,\mathbb Q)
   = (\mathbb Q a\oplus\mathbb Q b)/\mathbb Q(a+b).
\]
The same description follows from the etale Gysin sequence, with
the Tate twists understood.  For $m=1$ use the reduced quotient
$H^0(Q,\mathbb Q)/\mathbb Q\cdot1$, which is again the displayed
one-dimensional quotient.  Restriction to $Q$ and this quotient
are the local defect map, so $ua+vb$ maps to $u-v$, up to a
nonzero choice of orientation.  This proves the required
identification rather than assuming a local algebraic generator.
If
$\delta_s\ne0$, the indicated map is injective and therefore
$\nu_s$ is zero on all vertical cycles.  The conjugacy of the local
vanishing cycles (SGA 7 II, Expose XVIII, 6.6; Deligne, \emph{Weil I},
(5.4) and (5.8)) shows that this applies to every singular fibre
when $E_{\mathrm{van}}\ne0$.

If $E_{\mathrm{van}}=0$, the fixed-part theorem for the pencil gives
an isomorphism of rational Hodge structures and of variations
\[
 H^{2m-1}(X,\mathbb Q)|_U\xrightarrow{\sim}
 R^{2m-1}f_*\mathbb Q|_U,
\]
where $U$ is the smooth locus of the base.  Here injectivity is weak
Lefschetz, and surjectivity is the invariant-cycle statement:
the vanishing part is the orthogonal complement of the fixed part.
Equivalently one can use Deligne, \emph{Weil I}, (5.8)(b)(1), together
with the global invariant-cycle theorem.  This makes the variation
constant as a variation of Hodge structure, not just as a local
system.

The $H^1$ of the higher Picard variety $P=\operatorname{Pic}^m_Y$,
with its Tate twist, is a direct summand of this odd cohomology.
Here are the maps.  Let $g=\dim P$; the case $g=0$ is immediate.
A Poincare cycle $\Gamma\in\operatorname{CH}^{m}(P\times Y)_{\mathbb Q}$
and its transpose induce
\[
 u:H^{2g-1}(P)\longrightarrow H^{2m-1}(Y),\qquad
 v:H^{2m-1}(Y)\longrightarrow H^1(P),
\]
with the required Tate twists understood.  The divisor
correspondence giving $vu$ induces a nonzero multiple of
$\lambda_Y^m:P\to P^\vee$, since $\Gamma$ induces a nonzero
multiple of the identity on the higher Picard variety.  Hence
$vu$ is invertible.  Thus $u$ is split injective, and the
Lefschetz isomorphism on the abelian variety $P$ identifies its
source with a twist of $H^1(P)$.  The cycle may be chosen over
$\mathbb C(t)$ with rational coefficients: take the trace of a
Poincare cycle defined over a finite Galois extension, which only
changes its nonzero multiplicity.  The construction therefore
gives the required maps of variations over a nonempty open $U$.
The existence and isogeny statements here are Saito,
\emph{Abelian varieties attached to cycles of intermediate
dimension}, Section 4, or Chen, Lemmas 3.16 and 3.19.
Thus the polarized weight-one variation of $P$ is constant, and
$P/\mathbb C(t)$ is isogenous to a constant abelian variety.
Indeed a subvariation of a constant polarizable Hodge structure
is specified by a fixed rational subspace, whose Hodge filtration
is its intersection with the fixed filtration; polarized weight-one
Hodge structures recover abelian varieties up to isogeny.

Every $\mathbb C(t)$-point of a constant abelian variety is constant:
the associated rational map from $\mathbb P^1$ extends by properness,
and its differential is zero because invariant differentials on
an abelian variety span its cotangent space whereas
$H^0(\mathbb P^1,\Omega^1)=0$.  The same assertion modulo a finite
group holds after isogeny.  Consequently the Poincare Neron--Tate
pairing of any two $\mathbb C(t)$-points of $P$ is zero.  Formula
(44.2) gives $\langle z,z'\rangle_G=0$, proving (44.11).
\end{proof}

\paragraph{What ``Cartier'' does and does not mean here.}
For $m=2$, vertical middle cycles are Weil divisors on threefold
fibres.  The zero-defect classes may be replaced, numerically for
this calculation, by the radical; this suffices.  The stronger
claim that \emph{every} such Weil divisor is Cartier is unnecessary
and is not asserted.  For $m>2$, middle vertical cycles are not
divisors at all.  The precise invariant needed by Theorem 5 is
the vertical intersection form, not a factoriality assertion.

\subsection{SN-44-04: three nodal fibres, three individual units}
\label{sn:44-04}

Use the pencil and notation of HODGE\_SIGN AR-01:
\[
 X=\mathbb P^2_x\times\mathbb P^2_y,\qquad
 Z=a^2-ab+b^2,\qquad
 M=\left\{\sum_{j=0}^2 d_j(\lambda,\mu)x_jy_j=0\right\}
       \subset X\times\mathbb P^1,
\]
with $(d_0,d_1,d_2)=(\lambda,\lambda+\mu,\lambda+2\mu)$.
The incidence model $M$ is the blowup of the smooth base locus.
Write $t=f^*[\mathrm{point}]$.  Its three singular parameters are
$s_0=[0:1]$, $s_1=[-1:1]$, and $s_2=[-2:1]$.  For each parameter
define the \emph{actual} surface
\[
 D_i=\{x=e_i,\ s=s_i\}\times\mathbb P^2_y\ \subset M.
 \tag{44.12}
\]
Then
\begin{gather}
 D_iD_j=0\ (i\ne j),\qquad \deg D_i^2=1,\qquad
 \deg(\pi^*Z\,D_i)=1,                                    \tag{44.13}\\
 \sum_{i=0}^2[D_i]=\pi^*Z+t(2a-b)
          \quad\hbox{in }\operatorname{CH}^2(M),           \tag{44.14}\\
 (\pi^*Z)^h=\pi^*Z-\sum_iD_i=t(b-2a).                    \tag{44.15}
\end{gather}
The last class is vertical and orthogonal to every vertical class.
Each nodal fibre therefore contributes precisely $1$ to the
subtraction in (44.10), and the corrected formula is
\[
 \boxed{\qquad 0=3-(1+1+1).\qquad}                       \tag{44.16}
\]

\begin{proof}
Let $Q_y$ be the rank-two universal quotient bundle on
$\mathbb P^2_y$, so that
\[
 0\longrightarrow\mathcal O(-b)\longrightarrow
 \mathcal O^{\oplus3}\longrightarrow Q_y\longrightarrow0.
\]
On $M$, the linear functional
$(d_0x_0,d_1x_1,d_2x_2)$ vanishes on the tautological line represented
by $y$, by the defining equation of $M$.  It therefore defines a
section of the rank-two bundle
\[
 \mathcal F=Q_y^\vee\otimes\mathcal O_M(a+t).
\]
Its zero scheme is exactly $D_0\sqcup D_1\sqcup D_2$.
To check the scheme structure, around $(e_i,s_i)$ on
$\mathbb P^2_x\times\mathbb P^1$, the three functions $d_jx_j$
are regular parameters: the two $d_j(s_i)$ with $j\ne i$ and
the derivative of $d_i$ are nonzero.  In $M$, their relation with
$y$ eliminates one of them on each standard $y$-chart.  The
remaining two are regular parameters normal to $D_i$.
Thus the section is regular and its zero scheme is reduced.

The top Chern class formula for a regular section gives
\[
 \begin{split}
 \sum_i[D_i]=c_2(\mathcal F)
 &=c_2(Q_y^\vee)+(a+t)c_1(Q_y^\vee)+(a+t)^2\\
 &=b^2-(a+t)b+(a+t)^2\\
 &=a^2-ab+b^2+t(2a-b),
 \end{split}
\]
using $t^2=0$.  This proves the integral Chow identity (44.14).
Since $a$ and $t$ restrict trivially to $D_i$, its normal bundle is
\[
 N_{D_i/M}=\mathcal F|_{D_i}=Q_y^\vee
             \simeq\Omega^1_{\mathbb P^2}(1).
\]
The Euler sequence gives $c(Q_y^\vee)=1-b+b^2$, and hence
$\deg D_i^2=\int_{\mathbb P^2}b^2=1$.  Distinct $D_i$ lie on
distinct fibres and are disjoint.  Finally
$\pi^*Z|_{D_i}=b^2$, proving all of (44.13).

For any vertical cycle $v$, projection to a point of the base gives
$t|_{\operatorname{supp}v}=0$ in the operational Chow ring.  Thus
$t(b-2a)$ has intersection zero with every vertical class; it is an
admissible extension of the zero generic class.  Its square is zero
because $t^2=0$.  Consequently the vertical projection of $\pi^*Z$
has components exactly $D_i$, and their three squares are the three
individual correction terms.  The already proved AR-01 facts
$Z|_Y=0$ and $\deg Z^2=3$ now give (44.16).
\end{proof}

This proof also identifies the local normalization in SN-44-02:
the strict transform of $D_i$ meets the exceptional quadric in one
ruling line, so $\nu_{s_i}(D_i)=1$ after choosing that ruling first.
The normal bundle computation independently verifies its norm.

\subsection{SN-44-05: the correction as a Tate-motive projector}
\label{sn:44-05}

\paragraph{Statement and proof.}
For a normalized nodal class $D_s$ in SN-44-02, put
\[
 p_s=(-1)^mD_s\times D_s\in\operatorname{CH}^{2m}(M\times M)_{\mathbb Q}.
\]
Composition of product correspondences multiplies their outside
factors by the middle intersection number.  Since
$D_s^2=(-1)^m$ and different fibres have zero intersection,
$p_s^2=p_s$ and $p_sp_r=0$ for $s\ne r$.
Each projector factors through, and has image isomorphic to,
the Tate motive $\mathbb L^m$: the two morphisms are supplied by
$D_s$ and $(-1)^mD_s$, whose composite on $\mathbb L^m$ is $1$.
Its action on a middle cycle class is
$p_s\alpha=(-1)^m\deg(\alpha D_s)D_s$.  Thus the corrected
extension in (44.10) is precisely the image of $\alpha$ under
$1-\sum_sp_s$, modulo the vertical radical.  This proves the
motivic assertion by explicit algebraic correspondences; it does
not assume a general standard conjecture. \qed

\paragraph{Schemes over $\mathbb Z$, $\lambda$-rings and $K$-theory.}
The example, all three $D_i$, the regular section of $\mathcal F$,
and its integral Chow identity are defined over
$\operatorname{Spec}\mathbb Z[1/2]$.  Over this base the three
parameters remain distinct, the coefficient differences needed
above are units, and the same local parameter check proves
relative smoothness of the total space and regularity of the
three zero components.  In $K_0(M)$ the Koszul resolution reads
\[
 [\mathcal O_{\coprod D_i}]
     =\lambda_{-1}(\mathcal F^\vee)
     =1-[\mathcal F^\vee]+[\det\mathcal F^\vee].
\]
Its degree-two Chern character is
$\operatorname{ch}_2[\mathcal O_{\coprod D_i}]
 =c_2(\mathcal F)=\pi^*Z+t(2a-b)$.
These identities are the $K$-theoretic and $\lambda$-ring versions
of (44.14).  Regular immersions and the normal-bundle sequence
belong to the framework of EGA IV$_4$, 19.1.1;
the intersection operations used here are the refined Gysin and
self-intersection operations of Fulton, Sections 6.2--6.3.

\paragraph{Etale sheaves, weights, motives and $L$-functions.}
For a fibre of the example over $\mathbb F_q$ of odd characteristic,
the classes $D_i$ give three algebraic Tate summands
$\mathbb Q_\ell(-2)\subset H^4(M_{\overline{\mathbb F}_q},\mathbb Q_\ell)$.
Their Frobenius eigenvalues are $q^2$, so their weights are $4$.
At sheaf level these are the three skyscraper defects in
$R^4f_*\mathbb Q_\ell$ allowed by Deligne, \emph{Weil I},
(5.8)(b)(2), in the case $E_{\mathrm{van}}=0$; the local
specialization sequence is (4.3.3), and the SGA locus is SGA 7 II,
Expose XIV, 3.2, and Expose XVIII, 6.6.  Their degree-four zeta
factor is $(1-q^2T)^{-3}$, and the corresponding motive is
$\mathbb L^2\oplus\mathbb L^2\oplus\mathbb L^2$.
Purity assigns these classes the expected weight but cannot delete
their intersection contribution.  This is the precise reason a
sign or weight calculation cannot repair the omitted term in
cote 44 by itself.  These are geometric height assertions over a
curve and fibrewise statements over $\mathbb Z[1/2]$; no assertion
about an Arakelov height over $\operatorname{Spec}\mathbb Z$ is made.

\subsection{SN-44-06: extension of the height comparison and the
\texorpdfstring{$E=0$}{E=0} conclusion to positive characteristic}
\label{sn:44-06}

This subsection supplies a separate algebraic proof of the extension;
it does not extend Chen's complex theorem by a change of notation.
The local calculation SN-44-02 and the explicit example already work
in characteristic different from $2$.  We now show that the
Neron--Tate interpretation and the zero-height assertion of SN-44-03
also hold in that range.  The height comparison itself works in every
characteristic.

\paragraph{Finite-extension normalization and descent.}
Let $k$ be algebraically closed, $B/k$ a smooth projective connected
curve, and $K=k(B)$.  An incidence Picard variety and its Poincare
cycle, incidence duality, and the finitely many points being paired
may first be defined over a finite extension $L/K$.  Compute their
usual geometric Poincare Neron--Tate pairing over the normalization
$B_L$ of $B$, and divide its degree by $[L:K]$.  This gives the
pairing relative to the original base $B$.  All finite extensions,
including inseparable ones, are normalized by their full degree.
The Poincare pairing and Kahn's geometric height have the same
base-change rule.  For the latter this is Proposition 3.9(iv),
extended to arbitrary regular generic fibres by Theorem 4.14(b)
of \cite{KahnHeight2024}; thus no resolution of the base-changed
high-dimensional model is needed.  The normalized number is
independent of the finite field of definition: compare over a
common finite extension.  It is an ordinary Neron--Tate pairing
on an abelian variety with this usual extension-of-fields
normalization.

For the $E=0$ assertion we need to preserve the rational base curve,
so descent only to an arbitrary finite separable extension would
not suffice.  Here is the more precise descent that does suffice.
Over any perfect field $F$, the incidence Picard variety of
$Y_{\overline F}$ descends to $F$.  Indeed the universal property
for incidence homomorphisms supplies, for each
$\sigma\in\operatorname{Gal}(\overline F/F)$, a unique compatible
isomorphism from its $\sigma$-twist, and uniqueness supplies the
cocycle identity.  This datum is continuous: the variety and a
correspondence realizing its incidence-universal map are defined
over a finite extension, and its stabilizer acts by their specified
base change.  Galois descent for abelian varieties is effective.
The same argument makes the incidence map equivariant, so a cycle
defined over $F$ has an $F$-rational image.  This is the
perfect-field descent recorded in Chen, Remark 3.21; Saito's
incidence construction and duality, Theorem 3.6 and Proposition 4.4,
apply over algebraically closed fields of arbitrary characteristic
\cite{SaitoPicard1979}.

In characteristic $p>0$, apply this to
$F=K^{\mathrm{perf}}$.  Trace a Poincare cycle from a finite Galois
extension to obtain a rational Poincare cycle over $F$; its induced
endomorphism is a positive integral multiple of the identity.
The abelian variety, this cycle, the duality, and the finitely many
points and identities all descend, by finite presentation, to some
finite purely inseparable subextension $K_r/K$ of $F/K$.
When $K=k(t)$ we can enlarge it to
$K_r=k(t^{1/p^r})$; its smooth projective base is again
$\mathbb P^1_k$.  This argument avoids asserting an unproved
descent statement over the imperfect field $k(t)$ itself.

\paragraph{The arbitrary-characteristic height comparison.}
With the finite-extension normalization just specified, (44.2)
holds over $K=k(B)$ in every characteristic for algebraically
trivial cycles.  Here is the reduction to the classical curve
pairing in full.  After a finite extension, algebraic equivalence
expresses $z$ as a finite sum of terms $\Gamma_*\gamma$, with
$C$ a smooth projective geometrically connected curve,
$\gamma=[c_0]-[c_1]\in\operatorname{CH}^1(C)$,
$c_0,c_1\in C(K)$, and
$\Gamma\in\operatorname{CH}^m(C\times Y)_{\mathbb Q}$.
It is enough to prove the identity for one such term, by
bilinearity.  This is the curve reduction in Kahn's proof of
Theorem 5.6(a), also used in Chen's proof of Theorem 7.9.

Admissibility is preserved by these correspondences, by Kahn,
Proposition 3.6 (and Proposition 4.12(c) after alterations).
The projection formula gives the adjunction
\[
 h_Y(\Gamma_*\gamma,w)=h_C(\gamma,\Gamma^*w).               \tag{44.17}
\]
To make the degree shift explicit, take the closure of $\Gamma$
in the fibre product of regular models and admissible lifts of
$\gamma,w$.  Both sides of (44.17) are the pushforward to $B$ of
the product of this closure with the two pulled-back lifts,
formed by refined Gysin maps.  Commutativity of these operations
and the projection formula identify them.  This is exactly the
calculation in Kahn's proof of Lemma 3.3, with the complementary
codimensions shifted by the degree of $\Gamma$; it is not a use
of the degree-zero-only statement of Proposition 3.8 at a wrong
degree.  If a regular model is unavailable, pull back to admissible
alterations, apply this calculation, and divide by their degrees,
as in Theorem 4.14(b).

For $C$, the regular surface model exists in all characteristics.
At each place, a rational extension of a degree-zero divisor that
has degree zero on every fibre component computes Neron's local
pairing by intersection.  This is P\'epin,
\emph{Neron's pairing and relative algebraic equivalence},
Definition 2.1 and Theorem 2.3 \cite{PepinNeron2011}.
The condition on component degrees is precisely the admissibility
condition for this surface.  Summing the local identities is the
global Neron pairing, with normalized valuations and degree on $B$.
Equivalently, it is the degree of the canonical rational extension
of the Poincare biextension.  Under the Abel--Jacobi map its
duality is $\lambda_C^1$, the negative of the theta polarization:
the transpose of the universal divisor pulls back to the given
degree-zero divisor on $C$.  Functoriality of Neron's pairing
therefore gives, in precisely the normalization of (44.2),
\[
 h_C(\gamma,\eta)
 =\langle\operatorname{AJ}\gamma,
                 \lambda_C^1\operatorname{AJ}\eta\rangle_{\mathcal P}.
                                                               \tag{44.18}
\]
This uses the established pairing on curves, not a
higher-codimension height conjecture or a conjecture on component
group duality.

Let $a=\operatorname{Pic}(\Gamma):J(C)\to P$ and
$b=\operatorname{Pic}(\Gamma^{\mathsf t}):P\to J(C)$.
Saito's functorial incidence-duality identity, Proposition 4.4,
specializes, since $k_C^1=1$, to
\[
 k_Y^m\lambda_C^1 b=a^\vee\lambda_Y^m.                    \tag{44.19}
\]
The incidence maps satisfy
$\theta(\Gamma_*\gamma)=a\operatorname{AJ}\gamma$ and
$\operatorname{AJ}(\Gamma^*w)=b\theta w$.
Combining (44.17)--(44.19) and functoriality of the Poincare
biextension proves
\[
\begin{split}
h_Y(z,w)
 &=\langle\operatorname{AJ}\gamma,
                 \lambda_C^1 b\theta w\rangle_{\mathcal P}\\
 &=\frac1{k_Y^m}\langle\operatorname{AJ}\gamma,
                 a^\vee\lambda_Y^m\theta w\rangle_{\mathcal P}\\
 &=\frac1{k_Y^m}\langle\theta z,
                 \lambda_Y^m\theta w\rangle_{\mathcal P}.
\end{split}
\]
Finally divide by the finite extension degree.  This proves
(44.2) in the stated generality.  If the traced Poincare cycle
has multiplicity $N k_Y^m$ instead of $k_Y^m$, its duality is
$N\lambda_Y^m$, so the normalized expression is unchanged.

The same argument covers complementary, possibly unequal,
codimensions $i+j=\dim Y+1$.  Choose a Poincare $i$-cycle on
$\operatorname{Pic}^i_Y\times Y$ of multiplicity $k_Y^i$; its
transpose gives
$\lambda_Y^j:\operatorname{Pic}^j_Y\to
(\operatorname{Pic}^i_Y)^\vee$.  Saito's same Proposition 4.4
then gives
\[
 h_Y(z,w)=\frac1{k_Y^i}
       \langle\theta_i z,\lambda_Y^j\theta_j w\rangle_{\mathcal P}.
\]
Indeed the curve correspondence has codimension $i$, its transpose
sends codimension $j$ to divisors, and (44.19) uses $k_Y^i$.
This supplies the Neron--Tate comparison for the nonmiddle
vanishing assertion in SN-44-01 as well.

\paragraph{The fixed part over $\mathbb P^1$ in arbitrary characteristic.}
Let the hypotheses of SN-44-03 hold over an algebraically closed
field of characteristic different from $2$, with $\ell$-adic
cohomology in place of Betti cohomology, and suppose
$E_{\mathrm{van}}=0$.  Deligne, \emph{Weil I}, (5.8)(b)(1),
says that $R^{2m-1}f_*\mathbb Q_\ell$ is constant.  Thus
$H^{2m-1}(Y_{\overline K},\mathbb Q_\ell)$ is a trivial
Galois representation.  Use the descent above to work over
$K_r=k(t^{1/p^r})$ in characteristic $p$ (and $K_r=K$ in
characteristic zero), so that the incidence Picard variety $P$
and a rational Poincare cycle are defined there.

The maps $u,v$ in the proof of SN-44-03 have the same construction
in $\ell$-adic cohomology.  The isogeny $vu$ makes $v$
surjective.  In particular $H^1(P_{\overline K},\mathbb Q_\ell)$,
up to its Tate twist, is a quotient of the trivial middle
representation, so is itself trivial.  All the Tate twists are
trivial on this geometric Galois group because $k$ is algebraically
closed.  The integral Tate module is also unramified: an
endomorphism of its free $\mathbb Z_\ell$ lattice that is zero
after tensoring with $\mathbb Q_\ell$ is zero.

The Neron--Ogg--Shafarevich criterion (SGA 7 I, Expose IX,
Corollary 2.2.9) now gives good reduction of $P$ at every point
of $\mathbb P^1_k$.  Its Neron model
$g:\mathcal A\to\mathbb P^1_k$ is therefore an abelian scheme.
It is projective: the global Neron model over a noetherian regular
curve is quasi-projective (SGA 7 I, IX, Proposition 3.2 and
3.2.1, printed pp.~347--348, citing Raynaud, VIII.2(3))
\cite{SGA7IHeight}, and good reduction makes it
proper.  Its total space is consequently smooth projective over
$k$.  We spell out the fixed-part argument to avoid confusing
constant up to isogeny with constant as a family.

The low-degree Leray exact sequence for $g$ is
\[
0\longrightarrow H^1(\mathbb P^1,\mathbb Q_\ell)
 \longrightarrow H^1(\mathcal A,\mathbb Q_\ell)
 \longrightarrow H^0(\mathbb P^1,R^1g_*\mathbb Q_\ell)
 \xrightarrow{d_2}H^2(\mathbb P^1,\mathbb Q_\ell)
 \longrightarrow H^2(\mathcal A,\mathbb Q_\ell).
\]
The first term is zero.  The last arrow is injective because the
zero section is a left inverse to $g^*$, so $d_2=0$.
Moreover $R^1g_*\mathbb Q_\ell$ is constant, as is already true
at the generic point.  Hence restriction gives
\[
 H^1(\mathcal A,\mathbb Q_\ell)
   \xrightarrow{\sim}H^1(\mathcal A_s,\mathbb Q_\ell)
 \quad\text{for every }s\in\mathbb P^1(k).                 \tag{44.20}
\]

Put $A_0=\operatorname{Pic}^0_{\mathcal A/k,\mathrm{red}}$,
an abelian variety over $k$.  A Poincare bundle exists after
rigidification at a $k$-point of $\mathcal A$, and restriction
to the fibres defines a homomorphism of abelian schemes
\[
 A_0\times_k\mathbb P^1\longrightarrow\mathcal A^\vee.
\]
By the Kummer identification
$V_\ell(\operatorname{Pic}^0(T)_{\mathrm{red}})
 =H^1(T,\mathbb Q_\ell(1))$ for smooth projective $T/k$,
(44.20) says that its maps on rational Tate modules are
isomorphisms.  In particular its generic fibre is an isogeny.
Dualizing, $P$ is isogenous over $K_r$ to a constant abelian
variety.  This conclusion permits nonconstant abelian schemes
that are constant only up to isogeny.

Every $K_r$-point of a constant abelian variety is constant:
the rational map from $\mathbb P^1_k$ extends by properness,
and every morphism $\mathbb P^1_k\to A_0$ is constant, by the
Albanese universal property and
$\operatorname{Alb}(\mathbb P^1)=0$.  This argument works also
for inseparable morphisms.  A quasi-inverse to the isogeny shows
that $P(K_r)$ modulo its $K_r/k$-trace is torsion.  The canonical
Poincare Neron--Tate pairing kills the trace and torsion, so it
vanishes on the points coming from $z,z'$.  The comparison proved
above and division by $[K_r:K]$ now give
$h_{\mathrm{geom}}(z,z')=0$ over the original base.

The $E\ne0$ part of SN-44-03 also carries over: use the etale
local specialization sequence (Deligne, (4.3.3)) and the
conjugacy in (5.8); the difference-of-rulings defect is the
algebraic realization of the same one-dimensional local term.
Thus (44.10) and (44.11), including their Neron--Tate
interpretation with the specified finite-extension normalization,
hold for Lefschetz pencils of odd-dimensional members at least
one over algebraically closed fields of characteristic
different from $2$.  The general Gram formula (44.4) remains
valid in characteristic $2$ whenever its stated smooth model
exists; the ordinary-node calculation used here is not asserted
in that characteristic. \qed
